% Paper 3: computer-assisted proof that the known 9-point Heilbronn configuration in the unit square is optimal.
% Every numerical claim carries a hidden "% SRC:" comment naming the lab file it comes from.
\documentclass[11pt]{amsart}

\usepackage[T1]{fontenc}
\usepackage[utf8]{inputenc}
\usepackage{lmodern}
\usepackage{microtype}
\usepackage[margin=2.6cm]{geometry}
\usepackage{amsmath,amssymb,amsthm}
\usepackage{booktabs}
\usepackage{array}
\usepackage{xcolor}
\usepackage{enumitem}
\usepackage{url}
\usepackage[colorlinks=true,linkcolor=blue!50!black,urlcolor=blue!50!black,citecolor=blue!50!black]{hyperref}

% Public repository accompanying the paper (single definition).
\newcommand{\repourl}{\url{https://github.com/NERD-NEH007/heilbronn-discovery-lab}}

\setlist{itemsep=1pt,topsep=3pt}

\newtheorem{theorem}{Theorem}[section]
\newtheorem{proposition}[theorem]{Proposition}
\newtheorem{lemma}[theorem]{Lemma}
\newtheorem{corollary}[theorem]{Corollary}
\theoremstyle{definition}
\newtheorem{definition}[theorem]{Definition}
\newtheorem{remark}[theorem]{Remark}

\newcommand{\Q}{\mathbb{Q}}
\newcommand{\R}{\mathbb{R}}
\newcommand{\sq}{[0,1]^2}
\newcommand{\tstar}{t^{*}}

\begin{document}

\title[The 9-point Heilbronn configuration is optimal]{A computer-assisted proof that the known nine-point Heilbronn configuration in the unit square is optimal}

\author{Ethan Hendren}
\email{eehendren007@gmail.com}

\begin{abstract}
For Heilbronn's problem in the unit square we give a computer-assisted proof that the largest possible minimum triangle area of nine points is $H(9) = (9\sqrt{65}-55)/320 = 0.0548759991\ldots$, and that it is attained only by the known configuration (listed on the record site under Comellas and Yebra) and its images under the symmetries of the square.
Optimality for $n = 9$ was previously established by global optimization software within numerical tolerances, not rigorously.
The proof is a branch and bound over boxes of point positions on a dyadic grid, after exact symmetry and normalisation reductions. Every discarded box is discarded by a test evaluated in exact integer or rational arithmetic; a floating-point LP solver only proposes dual multipliers, which are rounded and re-checked exactly. The boxes that survive are covered by a separately computed exact local certificate around the optimum. The run took $3.59\cdot10^{8}$ nodes and 16 hours on 14 worker processes, with no failures.
\textbf{This is an AI-assisted, computer-assisted proof. It has not been independently checked}: the code was written with an AI system, it is not formally verified, and nobody other than the author has re-run or reviewed it. Section~\ref{sec:limits} states what an independent check would require.
\end{abstract}

\maketitle

% =====================================================================
\section{Introduction}\label{sec:intro}

Let $Q = \sq$ and, for $P = (p_1,\dots,p_n) \in Q^n$, let
\[
  f(P) = \min_{i<j<k} \lvert\Delta(p_i,p_j,p_k)\rvert, \qquad \Delta(a,b,c) = \tfrac12 \det(b-a,\, c-a),
\]
the smallest area of a triangle with vertices in $P$. Heilbronn's problem asks for $H(n) = \max_{P \in Q^n} f(P)$ and for the maximising configurations. The maximum exists because $f$ is continuous and $Q^n$ is compact.

\subsection*{Prior work}
Optimality is proved for $n \le 8$: $n = 5$ and $n = 6$ by Yang, Zhang and Zeng \cite{YZZ,YZZ92} ($n = 6$ in \cite[Theorem~3]{YZZ92}; first solved by Yang and Zhang \cite{YZ80}; see also Dress, Yang and Zeng \cite{DYZ}), $n = 7$ by Zeng and Chen \cite{ZC}, $n = 8$ by Dehbi and Zeng \cite{DZ}; these are the sources listed by the record site \cite{Stead}. % SRC: paper1/heilbronn_structure.tex l.73; papers/citation_findings.md Sec. 1 (record site)
The nine-point configuration studied here was found by Comellas and Yebra \cite{CY}; the record site lists it under their names \cite{Stead}. % SRC: papers/citation_findings.md Sec. 1 (record site holders: n = 7, 8, 9, 10, 12 Comellas & Yebra, Dec 2001)
For $n = 9$, Chen, Xu and Zeng \cite{Chen17} confined all optimal configurations by a GPGPU search and obtained $0.054875999 < H(9) < 0.054878314$ (as quoted in \cite{MMK}); this bounds $H(9)$ but does not prove the optimality of the known configuration. % SRC: paper1 l.76; citation_findings.md Sec. 1 (MMK / Chen et al. 2017)
Monji, Modir and Kocuk \cite{MMK} give MIQCP formulations with symmetry breaking and solve $n \le 9$. % SRC: paper1 l.75
Sudermann-Merx \cite{SM} solves $n \le 10$ by mixed-integer nonlinear programming and states that $n = 9$ is solved to global optimality within the numerical tolerances of the solver, that proofs exist for $n \le 8$ but not for $n = 9$ and $10$, and gives the exact value $(9\sqrt{65}-55)/320$. % SRC: citation_findings.md Sec. 1 (SM, n = 9, 10; exact forms)
Both optimization results certify the known configuration within solver tolerances. The record site lists $n = 9$ as ``numerically optimal''. Panorias \cite{Pan10} has since deposited an exact computer-assisted proof of the value for $n = 10$ (13 November 2026; AI-assisted and not yet independently checked; it does not classify the maximizers). % SRC: citation_findings.md Sec. 1 (record site, n = 9 and 10)

\subsection*{Result}
Put
\[
 a = 1 - \tfrac{\sqrt{65}}{10},\quad b = \tfrac58 + \tfrac{\sqrt{65}}{40},\quad c = \tfrac{9}{16} - \tfrac{3\sqrt{65}}{80},\quad d = \tfrac34 - \tfrac{\sqrt{65}}{20},
\]
so $a \approx 0.19377$, $b \approx 0.82656$, $c \approx 0.26017$, $d \approx 0.34689$, and let
\[
 C_9 = \bigl( (a,1),\ (1,a),\ (1,b),\ (b,1),\ (0,b),\ (b,0),\ (c,0),\ (0,c),\ (d,d) \bigr).
\]
% SRC: papers/paper3/check_closed_form_exact.py (exact identification of the exp 44 / exp 103 certified centre in Q(sqrt 65))
$C_9$ is symmetric under the reflection in the diagonal $y = x$; exactly $11$ of its $84$ triangles have the minimum area, and every other triangle is strictly larger, the smallest excess being $\approx 1.529\cdot10^{-3}$ (Lemma~\ref{lem:ident}).

\begin{theorem}\label{thm:main}
Let $\tstar = (9\sqrt{65}-55)/320 = 0.054875999170896\ldots$. Every configuration of nine points in the unit square has a triangle of area at most $\tstar$, that is, $H(9) = \tstar$. A configuration $P \in Q^9$ has $f(P) = \tstar$ if and only if, after relabelling, $P = g(C_9)$ for one of the eight symmetries $g$ of the square; these are four distinct point sets.
\end{theorem}
% SRC: 110-n9-proof/NOTES.md "Result (2026-11-13 12:31)"; results/run_n9/stats.json

The value $\tstar$ and the configuration are not new; what is new is a proof. The proof is computer-assisted (Sections~\ref{sec:outline}--\ref{sec:run}).

\begin{remark}[Status of this proof]\label{rem:status}
This is an \emph{AI-assisted, computer-assisted proof that has not been independently checked.} The branch-and-bound code, the local certificate and the reductions were developed with the help of an AI system (see ``Use of AI tools''). The code is not formally verified, the reductions of Section~\ref{sec:reductions} have been checked only by the author, and the 16-hour computation has been run once, on one machine. We have tested the implementation (Section~\ref{sec:validation}), but tests are not a check of the proof. Until an independent check of the kind described in Section~\ref{sec:limits} has been carried out, Theorem~\ref{thm:main} should be read as a claim supported by a computation, not as an established result.
\end{remark}

% =====================================================================
\section{Outline of the proof}\label{sec:outline}

The proof has three parts.
\begin{enumerate}[label=(\Alph*)]
\item \emph{A local certificate} (Theorem~\ref{thm:local}): an exact configuration $x^*$ with value $\tstar$, and a sup-norm ball around it in which every other configuration has minimum area $< \tstar$.
\item \emph{A global exclusion} (Proposition~\ref{prop:bb}): after exact normalisations, every configuration with minimum area above a threshold $T$ slightly below $\tstar$ lies inside that ball, up to a symmetry of the square and a relabelling. This is the branch-and-bound computation.
\item \emph{The identification} $x^* = C_9$ and $f(x^*) = (9\sqrt{65}-55)/320$ (Lemma~\ref{lem:ident}).
\end{enumerate}

\subsection{The local certificate}
The local certifier is the one of our experiment~44, run unchanged on a numerical $n = 9$ optimum; we summarise what it proves. It works with the active system: the triangles whose area equals the minimum (``ties'') and the wall contacts. Using the diagonal symmetry of the configuration it parametrises symmetric, wall-respecting configurations affinely by parameters $u$, solves the tie equations by Newton's method on a grid of $2^{-110}$ with exact residuals, and proves by a Krawczyk (radii-polynomial) test in exact rational arithmetic that the tie system has exactly one zero $(u^*,\tstar)$ in the sup-ball of radius $r$ around the polished point $(\tilde u, \tilde t)$. % SRC: 44-interval-certificate/NOTES.md steps 1-4; 103-global-proof/get_certs.py
For $n = 9$: no rattlers (all nine points are in the core), $11$ loaded ties and $8$ loaded wall contacts, $19 = 2\cdot 9+1$ constraints, first-order mode, $r = 4.65\cdot10^{-32}$, and every point moves by at most $\rho = r$ (the affine map $u \mapsto x$ has row sums of absolute values at most~$1$). The value is enclosed in $[t_{\rm lower}, t_{\rm upper}]$ with
$t_{\rm lower} = 0.05487599917089670897281099710224$ and $t_{\rm upper} = t_{\rm lower} + 10^{-31}$. % SRC: 103-global-proof/results/local_certs.json, key "9": mode first-order, n_core 9, rattlers [], loaded_ties 11, loaded_walls 8, existence_radius = box_rho = 4.649e-32, t_lower, t_upper

The certifier then proves that every triangle that is not exactly tied stays strictly above $\tstar$, and that the $19$ loaded constraints carry strictly positive multipliers (lower bound $0.0217$). It uses these as follows. % SRC: local_certs.json "9": y_min_lower 0.0217; 44/NOTES.md steps 6-7

\begin{theorem}[local certificate; experiment 44]\label{thm:local}
There is a configuration $x^* \in Q^9$, with $\|x^* - \tilde x\|_\infty \le 4.65\cdot 10^{-32}$ for an explicit rational configuration $\tilde x$, such that $f(x^*) = \tstar \in [t_{\rm lower}, t_{\rm upper}]$ and
\[
  f(x) < \tstar \quad\text{for every } x \in Q^9 \text{ with } 0 < \|x - x^*\|_\infty < \rho_b, \qquad \rho_b \ge 1.1786\cdot10^{-4}.
\]
\end{theorem}
% SRC: local_certs.json "9": basin_radius 0.00011786328420762071; 44/NOTES.md step 7 "Theorem used"; certify_interval.py l.517-547

\begin{proof}[Sketch of the argument implemented by the certifier]
Let $M$ be the $19\times19$ matrix whose rows are $[\sigma_T\nabla\Delta_T,\,-1]$ for the $11$ ties ($\sigma_T$ the sign of the area) and $[n_w, 0]$ for the $8$ wall contacts ($n_w$ the inward normal), acting on $(\delta x, \tau)$. With an approximate inverse $C$, the quantity $\beta = \|I - CM(\tilde x)\|_\infty + 6\rho\|C\|_\infty$ is bounded exactly; $\beta < 1$ proves $M(x^*)$ nonsingular with $\|M^{-1}\|_\infty \le \|C\|_\infty/(1-\beta)$, and the entries of $y = -(M^{-1})_{\text{last},\cdot}$ are bounded below by $y_{\min} > 0$. For a displacement $\delta x$ that keeps the points in $Q$, with $\delta = \|\delta x\|_\infty$, put $\tau = \min_T \sigma_T\nabla\Delta_T\cdot\delta x$ and $s = M(\delta x,\tau) \ge 0$. Then $-\tau = y\cdot s \ge y_{\min}\|s\|_\infty \ge c\,\delta$ with $c = y_{\min}/\|M^{-1}\|_\infty$. The quadratic remainder of a triangle area is at most $4\delta^2$, so some tie has $\sigma_T\Delta_T(x^*+\delta x) \le \tstar - c\delta + 4\delta^2 < \tstar$ for $0 < \delta < c/4$; the radius is also capped at $\tilde t/10$ so that the signs $\sigma_T$ do not change. All bounds are computed in exact rational arithmetic; floats only choose $C$ and the active set. (One exception, not used in the basin argument: in the check that non-tied triangles exceed $\tstar$, triangles whose double-precision area exceeds $\tilde t$ by more than $10^{-6}$ are not re-checked exactly; for $n = 9$ this is superseded by the exact check of Lemma~\ref{lem:ident}.) The certifier stores $\rho_b = \min(c/4, \tilde t/10)$ rounded to a double. Appendix~\ref{app:basin} gives the argument in full, with all constants for $n = 9$. % SRC: 44/NOTES.md step 7; certify_interval.py l.535-545
\end{proof}

\subsection{The global exclusion}
Let $T = 494279058340989 / 2^{53} = 0.0548759991160207\ldots$, the double immediately below the double nearest to $t_{\rm lower}(1-10^{-9})$; so $T < t_{\rm lower} \le \tstar$ (checked exactly). Let % SRC: results/run_n9/stats.json "T_exact"; bb2.py make_problem (mode exact)
\[
  R = \rho_b'(1 - 10^{-6}) - \rho - 2^{-60},
\]
where $\rho_b' = 1.1786328420762071\cdot10^{-4}$ is the stored (rounded) basin radius. Since $\rho_b'$ is a double rounded from the exact bound, $\rho_b'(1-10^{-6}) < \rho_b$, and so $R + \rho < \rho_b$. % SRC: 103-global-proof/bb.py load_covers; NOTES.md (103) Limitations "basin radius is read from a float ... reduced by a relative 1e-6"

\begin{proposition}[computational]\label{prop:bb}
Let $\Omega \subset Q^9$ be the normalised domain of Lemma~\ref{lem:norm}. For every $P \in \Omega$ with $f(P) > T$ there are a symmetry $g$ of the square and a permutation $\pi$ of the labels with $\|P - \pi(g(\tilde x))\|_\infty \le R$.
\end{proposition}
Proposition~\ref{prop:bb} is what the branch and bound of Sections~\ref{sec:reductions}--\ref{sec:sound} establishes, by the computation of Section~\ref{sec:run}.

\subsection{The identification}
\begin{lemma}\label{lem:ident}
$x^* = C_9$ up to the labelling, and $\tstar = (9\sqrt{65}-55)/320$. Exactly $11$ triangles of $C_9$ have area $\tstar$, and all others have area $> \tstar$ (the smallest excess is $\approx 1.529\cdot10^{-3}$).
\end{lemma}
\begin{proof}
The script \texttt{check\_closed\_form\_exact.py} (in the repository) uses \texttt{sympy} only to \emph{propose} $C_9$ from the rational centre $\tilde x$. It then checks the following in exact arithmetic in $\Q(\sqrt{65})$, with numbers stored as pairs $(a,b)$ of rationals meaning $a + b\sqrt{65}$ and signs decided exactly (by comparing $a^2$ with $65b^2$); it stops on any failure. Put $v = (9\sqrt{65}-55)/320$.
(1) $t_{\rm lower} < v < t_{\rm upper}$.
(2) Every coordinate of $C_9$ lies in $[0,1]$, and the coordinates of $\tilde x$ equal to $0$ or $1$ (its eight wall contacts) are the same coordinates of $C_9$, with the same values.
(3) Every coordinate of $C_9$ is within $r(1-10^{-6})$ of the same coordinate of $\tilde x$, where $r$ is the stored existence radius (the margin covers the double rounding of $r$); the distance is $\approx 2.32\cdot10^{-32}$.
(4) For each of the $84$ triangles, $\lvert\Delta(C_9)\rvert$ is either exactly $v$ or strictly greater than $v$; exactly $11$ are equal to $v$.
(5) The $11$ tied triangles of $C_9$, with their orientation signs, are exactly the triangles of $\tilde x$ whose area is minimal to relative precision $10^{-20}$, with the same signs; and the point set $C_9$ is invariant under $(x,y)\mapsto(y,x)$.

From (3) and (5): the points of $\tilde x$ are at pairwise sup distance $> 0.17$, and each point of $C_9$ is within $10^{-31}$ of the point of $\tilde x$ with the same label, so the diagonal reflection permutes the labels of $C_9$ exactly as it permutes those of $\tilde x$. With (2), $C_9 = p(u_C)$ for the affine parametrisation $p$ of the certifier (Appendix~\ref{app:basin}, \eqref{eq:param}), in which every parameter is a coordinate of a point up to sign and an additive constant; hence $\|u_C - \tilde u\|_\infty \le \|C_9 - \tilde x\|_\infty < r$. By (4) and (5), $(u_C, v)$ is a zero of the reduced tie system. Its last component satisfies $\lvert v - \tilde t\rvert < 10^{-32} < r$: the certifier prints $\tilde t$ rounded down to $32$ decimals as $0.05487599917089670897281099710229$, and $v$ exceeds this by $3.6\cdot10^{-33}$. So $(u_C, v)$ lies in the Krawczyk ball of radius $r$ around $(\tilde u,\tilde t)$, and by uniqueness of the zero in that ball (Lemma~\ref{lem:krawczyk}), $C_9 = x^*$ and $\tstar = v$. By (4), $f(C_9) = v$.
\end{proof}
% SRC: papers/paper3/check_closed_form_exact.py output (ALL EXACT CHECKS PASSED; dist 2.32e-32, 11 tied, smallest excess 0.001529); re-run of certify_interval for n = 9: t_polished 5487599917089670897281099710229e-32, parametrisation printed from build_param

\begin{remark}\label{rem:ident}
Two inputs of this proof are not stored in \texttt{local\_certs.json} and are taken from a re-run of the (deterministic) certifier on the same input, which reproduces the stored record exactly ($r$, $\rho$, $\rho_b$, $t_{\rm lower}$, $t_{\rm upper}$, the points): the 32-digit display of $\tilde t$, and the parametrisation \eqref{eq:param}. The exact script checks only the coordinates of $C_9$, not the component $\tilde t$. The identification is not needed for the optimality statement itself: without it, the proof gives $H(9) = f(x^*) \in [t_{\rm lower}, t_{\rm upper}]$ with $x^*$ unique up to symmetry, and $C_9$ agrees with $x^*$ to $10^{-31}$.
\end{remark}

\subsection{Proof of Theorem~\ref{thm:main} from (A)--(C)}
Let $P \in Q^9$ with $f(P) > T$. By Lemma~\ref{lem:norm} there are stretch factors $\lambda_x,\lambda_y \ge 1$, a symmetry $h$ of the square and a relabelling such that the resulting $P' \in \Omega$ satisfies $f(P') = \lambda_x\lambda_y f(P) \ge f(P) > T$. By Proposition~\ref{prop:bb}, $\|P' - \pi(g(\tilde x))\|_\infty \le R$, hence $\|g^{-1}(\pi^{-1}P') - x^*\|_\infty \le R + \rho < \rho_b$ (symmetries of the square are isometries of the sup norm, and relabelling and symmetries preserve $f$). By Theorem~\ref{thm:local}, $f(P') \le \tstar$, with equality only if $\pi^{-1}P' = g(x^*)$. Hence $f(P) \le \tstar$ for every $P$ (if $f(P) \le T$ this holds trivially), and $f(x^*) = \tstar$, so $H(9) = \tstar$.

If $f(P) = \tstar$, then $\tstar \le f(P') \le \tstar$, so $\lambda_x\lambda_y = 1$: $P$ already touches all four sides, $P'$ is $h(P)$ relabelled, and $P' = g(x^*)$ up to labels. So $P$ is a symmetry image of $x^* = C_9$. Conversely every such image has $f = \tstar$. Since the stabiliser of $C_9$ in the symmetry group of the square is $\{1, \text{diagonal reflection}\}$, the images form four distinct point sets. \qed
% SRC: 103-global-proof/NOTES.md "Statement proved" (exact mode); the equality-case argument via the stretch factor is spelled out here (NOTES give it implicitly)

% =====================================================================
\section{Reductions}\label{sec:reductions}

Coordinates are written $x_i, y_i$; a configuration is a point of $Q^9 \subset \R^{18}$.

\begin{lemma}[normalisation]\label{lem:norm}
Let $P \in Q^n$ with $f(P) > 0$. There is a configuration $P'$ with $f(P') = \lambda f(P)$ for some $\lambda \ge 1$, obtained from $P$ by axis-parallel stretches with factors $\lambda_x, \lambda_y \ge 1$ ($\lambda = \lambda_x\lambda_y$), a symmetry of the square and a relabelling, such that $P'$ lies in
\[
\Omega = \Bigl\{\, P \in Q^n :\ \begin{aligned}[t]
 &x_0 \le x_1 \le \dots \le x_{n-1},\quad x_0 = 0,\quad x_{n-1} = 1,\quad y_b = 0,\ y_t = 1 \text{ for some } b < t,\\
 &\tfrac1n\textstyle\sum_i x_i \le \tfrac12,\quad \sum_i (x_i - \tfrac12)^2 \ge \sum_i (y_i - \tfrac12)^2 \,\Bigr\}.
\end{aligned}
\]
If $\lambda = 1$ then $P$ touches all four sides and $P'$ is the image of $P$ under a symmetry of the square, relabelled.
\end{lemma}

\begin{proof}
\emph{Stretching.} If $x_{\max} - x_{\min} < 1$, the map $x \mapsto (x - x_{\min})/(x_{\max}-x_{\min})$ keeps $P$ in $Q$ and multiplies every triangle area by $\lambda_x = 1/(x_{\max}-x_{\min}) > 1$; the same in $y$. ($f(P) > 0$ excludes $x_{\max} = x_{\min}$.) Afterwards some point lies on each of the four sides. % SRC: 103/NOTES.md Search domain item 2
\emph{Symmetries.} Every symmetry of the square preserves all triangle areas and the property of touching all four sides. Apply them in this order.
(i) If $Q_x := \sum(x_i - \frac12)^2 < Q_y := \sum(y_i-\frac12)^2$, apply the diagonal reflection $(x,y) \mapsto (y,x)$; it exchanges $Q_x$ and $Q_y$.
(ii) If $\frac1n\sum x_i > \frac12$, apply $x \mapsto 1-x$; it maps the mean to $1 - $ mean and preserves $Q_x$, $Q_y$ (each $(x_i - \frac12)^2$ is invariant) and all $y_i$.
(iii) Now label the points so that $x_0 \le \dots \le x_{n-1}$ (any order among equal $x$), and let $b$, $t$ be the smallest labels of points on the bottom and the top side; $b \ne t$ since no point is on both. If $b > t$, apply $y \mapsto 1-y$: it exchanges the bottom and top sides and so the roles of $b$ and $t$, and it preserves all $x_i$ (so the labelling and the mean of $x$) and $Q_x$, $Q_y$.
Each step preserves the conditions established before it, so $P'$ satisfies all of them; $x_0 = 0$ and $x_{n-1} = 1$ because the points touch the left and right sides and the labels are sorted by $x$. If $\lambda = 1$, no stretch was applied. % SRC: 110-n9-proof/NOTES.md (c) Lemma and proof; 103/NOTES.md Search domain items 1-3
\end{proof}

\begin{remark}\label{rem:normgap}
The notes of the computation state the lemma with ``finally relabel by $x$''; the labelling must be fixed before the $y$-mirror (iii), since $b$ and $t$ are defined through it. This is only a matter of order: (iii) does not change any $x_i$. We also note that the root boxes (Section~\ref{sec:tests}) impose only $y_b = 0$, $y_t = 1$ for a pair $b < t$, not the minimality of $b$ and $t$; this enlarges the searched set and costs only time.
\end{remark}

% =====================================================================
\section{The branch and bound}\label{sec:tests}

\subsection{Boxes}
Fix the grid $D = 2^{24}$. A \emph{box} is a product of closed intervals $[\ell_k/D, h_k/D]$, $k = 1,\dots,18$, with integer endpoints $0 \le \ell_k \le h_k \le D$; the points inside a box are arbitrary reals. There is one \emph{root box} for each pair $b < t$ ($36$ roots for $n = 9$): $x_0 = 0$, $x_8 = 1$, $y_b = 0$, $y_t = 1$, all other coordinates in $[0,1]$. The roots cover $\Omega$. % SRC: bb.py roots() domain 'basic'; progress.log start line "open 36"
A box is processed by the tests below in order; a box that is not discarded is split. The aim is to show: \emph{no configuration $P$ in $\Omega$ with $f(P) > T$ lies in a discarded box, unless the box is covered} (Section~\ref{sec:cover}).

\subsection{P0: propagation of the linear and diagonal constraints}
In exact integers at grid units: the ordering $x_i \le x_{i+1}$ and the mean condition $\sum x_i \le nD/2$ shrink the bounds (for example $h_i \le nD/2 - \sum_{j\ne i}\ell_j$). For the diagonal condition, with $a_i = 2x_i - D$ and $b_i = 2y_i - D$ in grid units, the condition is $\sum a_i^2 \ge \sum b_i^2$. Let $A_{\max} = \sum_i \max_{\rm box} a_i^2$ and $B_{\min} = \sum_i \min_{\rm box} b_i^2$. If $A_{\max} < B_{\min}$ the box contains no point of $\Omega$. Otherwise each $|b_i| \le \sqrt{A_{\max} - B_{\min} + \min b_i^2}$, and each $|a_i| \ge \sqrt{B_{\min} - A_{\max} + \max a_i^2}$ when the right side is positive; this may exclude one side of $x_i = \frac12$ or the whole box. All square roots are integer square roots, and every new bound is rounded outward, so that the shrunken box still contains every point of $\Omega$ that the old box contained. If some interval becomes empty the box is discarded. Points removed by P0 are not in $\Omega$, so P0 never removes a configuration of interest. % SRC: bb2.py k_propagate; 110/NOTES.md (c) "Use (exact integers, P0)"

\subsection{P1: exact corner ranges}
For a triple $(i,j,k)$, $2\Delta = (x_j - x_i)(y_k - y_i) - (x_k - x_i)(y_j - y_i)$ is affine in each of the three points separately, so its maximum and minimum over a product of rectangles are attained at one of the $4^3 = 64$ corner tuples. At grid units $2\Delta D^2$ is an integer of absolute value $< 2^{51}$ at corners, computed exactly in 64-bit integers and compared with $T_2 = \lfloor 2TD^2 \rfloor$. % SRC: bb2.py k_ranges; 110/NOTES.md (b)
If $\max |2\Delta D^2| \le T_2$ over the corners for some triple, every configuration in the box has $f \le T$, and the box is discarded. Otherwise, if $\min 2\Delta D^2 \ge -T_2$, every configuration in the box with $\lvert\Delta\rvert > T$ for this triple has $\Delta > T$; the orientation $\sigma = +1$ is then \emph{determined}. Symmetrically for $\sigma = -1$. (Since the corner values are integers, comparison with $\lfloor 2TD^2\rfloor$ is equivalent to comparison with $2TD^2$.)

\subsection{P2: McCormick LP with exactly verified duals}
Introduce $\tau$, the coordinates, and $w_{ab}$ for $a \ne b$ standing for $x_a y_b$. Every configuration in the box with $f > T$, together with $w_{ab} = x_a y_b$ and $\tau = \min_{\text{determined } T}\sigma_T\, 2\Delta_T$, satisfies:
(a) $\tau - \sigma_T\,2\Delta_T(w) \le 0$ for every determined triple ($2\Delta$ is linear in $w$);
(b) the four McCormick inequalities for each $w_{ab}$ on the box;
(c) the ordering rows $x_i - x_{i+1} \le 0$ and the mean row $\sum x_i \le n/2$;
and $\tau > 2T$. Write these rows, scaled by $D^2$ to integer coefficients, as $Gu \le H$. A floating-point LP solver (HiGHS) maximises $\tau$ and returns duals $\lambda \ge 0$. \emph{The proof does not use the LP value.} The duals of all rows of undetermined triples are set to zero; then $\mu = \lfloor \lambda 2^S\rfloor$ with $0 \le \mu < 2^{50}$, and $r = G^{\mathsf T}\mu$ is computed exactly (64-bit integer products of $G$ with the two 25-bit halves of $\mu$, recombined and summed in arbitrary-precision integers). Overflow is excluded statically: every entry of $G$ has absolute value at most $D = 2^{24}$ (the McCormick coefficients are box endpoints in $[0,D]$, the others are $\pm1$ or $\pm D$), each half of $\mu$ is below $2^{25}$, and no column of $G$ has more than $168$ nonzeros (the $\tau$ column, two rows per triple), so every 64-bit partial sum is below $168\cdot2^{49} < 2^{57}$; the code asserts the corresponding bound $<2^{62}$ when the model is built. This exact check is applied to \emph{every} box discarded by P2. With $s = r_\tau > 0$ the aggregated valid inequality $s\tau + \sum_{u\ne\tau} r_u u \le \mu^{\mathsf T}H$ gives
\[
\tau \le \frac{1}{s}\Bigl(\mu^{\mathsf T}H + \sum_{u\ne\tau}\max\bigl(-r_u \underline u,\, -r_u \overline u\bigr)\Bigr),
\]
where $[\underline u,\overline u]$ is the box interval of a coordinate or the exact range of $x_a y_b$ over the box. If the right side is $\le 2TD^2$ (checked as an integer inequality) the box is discarded. Rounding of the duals and solver inaccuracy can only make this test fail, never make it pass falsely. A box with $s \le 0$, or where the solver does not report an optimal solution, is simply split. Appendix~\ref{app:p2} lists every row and states and proves the certificate (Lemma~\ref{lem:farkas}). % SRC: bb2.py lp_bound2; 103/NOTES.md P2; 110/NOTES.md (a)

\subsection{Cover by the basin ball}\label{sec:cover}
If all box widths are at most $D/1000$, a float search proposes a symmetry $g$ and a matching $\pi$ of the box points to the points of $g(\tilde x)$; the box is then discarded if, in exact rational arithmetic, every endpoint of every coordinate interval is within $R$ of the matched coordinate of $g(\tilde x)$. Such a box lies in the sup-ball of radius $R$ around $\pi(g(\tilde x))$, which is exactly the conclusion of Proposition~\ref{prop:bb} for the configurations in it. % SRC: bb.py covered(); bb2.py make_problem cover_width=D//1000

\subsection{Splitting and failures}
Otherwise the box is split at the integer midpoint of a coordinate of width $\ge 2$ grid units (chosen by a heuristic: width times one plus the number of undetermined triples containing the point). The two closed children cover the parent. A box that is not discarded and has no coordinate of width $\ge 2$ is a \emph{failure}. The proposition is proved if the search, started from all root boxes, terminates with no failures. % SRC: bb2.py k_choose, process; 103/NOTES.md Splitting
An optional \emph{LP skip} omits P2 at a node whose parent's LP value exceeded $2T(1+1.2)$ (P0, P1 and the cover test are still applied); such a node is then split (or, at the grid floor, recorded as a failure), never discarded by P2, and its children always solve their LP. % SRC: 110/NOTES.md (d)

% =====================================================================
\section{Soundness}\label{sec:sound}

We list what has to be true for Proposition~\ref{prop:bb} to follow from a run with zero failures.
\begin{enumerate}[label=(S\arabic*)]
\item \emph{Coverage.} The root boxes contain $\Omega$ (Lemma~\ref{lem:norm} and Remark~\ref{rem:normgap}); P0 removes only points outside $\Omega$; splitting replaces a box by two closed boxes whose union is the box. Hence at every stage the union of the open boxes and the discarded boxes contains $\Omega$.
\item \emph{Discards.} Each discard is justified: P1 because some triple has $\lvert\Delta\rvert \le T$ on the whole box; P2 because the exact aggregated inequality bounds $\tau \le 2T$, contradicting $\tau > 2T$; the cover test by an exact containment in the $R$-ball. LP skips and solver failures never discard.
\item \emph{Exactness.} Every inequality used to discard is evaluated in integer or rational arithmetic: P0 and P1 in 64-bit integers with magnitudes below $2^{55}$ (no overflow is possible for $D = 2^{24}$: coordinates are below $2^{25}$, products below $2^{51}$, sums of nine squares below $2^{54}$); P2, for every P2 discard, in 64-bit integer products with partial sums below $2^{57}$ (static bound above; asserted $<2^{62}$) and arbitrary-precision sums; the cover test in rationals. Floating point only chooses the split coordinate, the LP duals, and the symmetry and matching in the cover test. A bad choice costs time, not soundness. % SRC: 103/NOTES.md "Why it is rigorous"; bb2.py k_matvec assertion; 110/NOTES.md (a), (b)
\item \emph{Termination.} The driver keeps a list of open boxes; a run is complete when this list and the set of boxes being processed are both empty. An exception in a worker (including an overflow assertion or an audit mismatch) propagates to the driver and stops the run without a completion record; the run completed in one session, with an empty error log. As an independent float tally, each root has weight $1/36$, a split divides a box's weight between its children in proportion to the two parts of the split interval, and the weights of closed boxes sum to $1$ up to rounding ($1 + 10^{-13}$). The exact check is the node count: the numbers of leaves and split nodes satisfy the identity of a forest of $36$ full binary trees (Section~\ref{sec:run}). Boxes at the grid floor that are neither discarded nor covered are counted as failures; there were none. % SRC: run_par2.py main loop; stats.json closed_w
\item \emph{Local theorem.} The cover is turned into Theorem~\ref{thm:main} by Theorem~\ref{thm:local}, whose own proof is the exact computation of experiment~44, with the margins $10^{-6}$ (relative) and $2^{-60}$ (absolute) absorbing the double rounding of the stored basin radius and the distance between $\tilde x$ and $x^*$.
\end{enumerate}

\noindent\emph{Trust base.} The branch-and-bound code (\texttt{bb.py} of experiment~103, about 500 lines, imported unchanged, and \texttt{bb2.py} of experiment~110, about 400 lines including numba-compiled kernels), the local certifier (about 600 lines), the Python integer and \texttt{Fraction} arithmetic, numba's 64-bit integer arithmetic (which wraps silently on overflow; overflow is excluded by the magnitude bounds above, not checked at run time), and the parallel driver's bookkeeping. HiGHS is outside the trust base. None of this is formally verified. % SRC: 103/NOTES.md Limitations; 44/NOTES.md Trust base; 110/NOTES.md Files

% =====================================================================
\section{Validation}\label{sec:validation}

These checks test the implementation; they are not part of the proof.
\begin{itemize}
\item \emph{Re-proofs for $n = 5$--$8$.} With the same code (experiment~110): $n = 5$ and $n = 6$ in the weaker form $H(n) \le \tstar_n(1+10^{-3})$, with node counts identical to those of experiment~103 ($1\,120$ and $9\,973$ without the diagonal reduction); $n = 7$ in exact mode (optimum unique up to symmetry) in $0.85$--$0.86$ million nodes, with and without the diagonal reduction and with LP skips; $n = 8$ in exact mode in $4\,252\,598$ nodes and $538$\,s on $14$ workers (experiment~103: $4\,256\,822$ nodes). All with zero failures. For $n = 5, 6$ only the weaker form is possible, since $n = 5$ is certified locally only at second order (no basin radius) and $n = 6$ has a one-parameter family of optima. The known optima for $n = 7, 8$ \cite{ZC,DZ} are reproduced. % SRC: 110/NOTES.md Validation table; 103/NOTES.md Results and "n = 5 and 6 are only eps-mode"
\item \emph{Random soundness test} (\texttt{test\_sound2.py}, $n = 7, 8, 9$, $600$ random boxes each near the optimum and its mirror images, with $T = 0.8\,\tstar_n$ so that the tests fire): no sampled configuration with all $\lvert\Delta\rvert > T$ lies in a box discarded by P1 or violates a determined orientation; no sampled point of the domain is removed by P0; no sample beats a certified P2 bound; the rational re-derivation agrees with the integer certificate on every tested box ($210$, $270$, $310$ LP boxes certified). A separate test of P0 on $20\,000$ random boxes at $n = 9$ ($5.6$ million domain samples) removed none. % SRC: 110/NOTES.md Validation bullets
\item \emph{Audit.} During the runs a random sample of P2 certificates (each already checked exactly by the integer code) is re-derived from the row definitions with rational arithmetic, independently of the integer code; a mismatch aborts. At $n = 5$ all $314$ were audited, at $n = 8$ $2\,701$; none disagreed. % SRC: 110/NOTES.md Validation table
\item \emph{Clean-room cross-check of P2.} A second checker (\texttt{clean\_room/check\_p2.py}) was written from Appendix~\ref{app:p2} and a neutral certificate format alone, without access to the search code. It recomputes the orientation pattern from the $64$ corner tuples of each box, rebuilds every row from its meaning, and decides each certificate in exact integer arithmetic. Certificates were dumped from the unmodified search code through its audit hook (\texttt{dump\_p2.py}). The checker accepted all $292\,120$ P2 certificates of a complete $n = 7$ run (whose node count, $853\,071$, equals that of the earlier run), a $10\%$ sample of $136\,972$ from a complete $n = 8$ run, and $8\,970$ sampled from four root boxes at $n = 9$: $438\,062$ certificates, none rejected. Its negative tests reject altered certificates (one multiplier halved, an area row's orientation flipped, the box widened, a row dropped, $T$ halved). For $n = 7$ this checks every P2 discard of that run independently; for $n = 9$ it is a sample, and a complete independent check would need per-leaf certificates from a full run. % SRC: 110/results/p2dump/*.log, clean_room/check_p2.py output 2026-11-13
\item \emph{Negative control.} With $n = 5$, $T = \tstar_5(1-10^{-3})$ and no cover, the claim is false, so the search must fail. It does: $5$ failures after $509$ nodes (with the diagonal reduction) and failures after $438$ nodes (without), at grid-floor boxes next to the optimum. % SRC: 110/NOTES.md "Negative control"; results/val/n5_negative_control*.txt
\end{itemize}

% =====================================================================
\section{The run}\label{sec:run}

The $n = 9$ run (exact mode, diagonal reduction on, LP skip threshold $1.2$, audit rate $0.02\%$) was started on 12 November 2026 and finished on 13 November 2026 without interruption, on a machine with 8 cores (16 threads) using 14 worker processes (HiGHS 1.16.0, numba 0.68.1). Table~\ref{tab:run} gives its statistics. % SRC: 110/RESUME.md command line; progress.log start and DONE lines; stdout (HiGHS version); 110/NOTES.md (numba version, machine)

\begin{table}[htbp]
\centering\small
\caption{Statistics of the $n = 9$ run.}\label{tab:run}
% SRC: 110-n9-proof/results/run_n9/stats.json
\begin{tabular}{@{}lr@{}}
\toprule
threshold $T$ & $494279058340989/2^{53}$ \\
cover radius $R$ & $\approx 1.1786\cdot10^{-4}$ \\
root boxes & $36$ \\
nodes & $358\,594\,738$ \\
\quad discarded by P0 & $2\,520\,863$ \\
\quad discarded by P1 & $70\,994\,112$ \\
\quad discarded by P2 & $105\,782\,380$ \\
\quad covered by the basin ball & $32$ \\
\quad failures & $0$ \\
LPs solved / LP skips & $278\,798\,572$ / $6\,281\,159$ \\
LP solver not optimal & $0$ \\
P2 certificates audited in rational arithmetic & $21\,286$ (all agreed) \\
maximum depth & $184$ \\
closed fraction of root volume & $1.0000000000000977$ \\
wall time & $57\,636$\,s ($16.0$\,h) \\
\bottomrule
\end{tabular}
\end{table}

Consistency: the discarded and covered boxes number $179\,297\,387$; the remaining $179\,297\,351$ nodes were split; a forest of binary trees with $36$ roots and $179\,297\,351$ internal nodes has $179\,297\,387$ leaves, as observed. Nodes reaching the LP stage ($358\,594\,738$ minus the P0, P1 and cover discards, $285\,079\,731$) equal LPs solved plus LP skips. The failure log is empty. The script \texttt{verify\_run.py} (in the repository) re-checks these facts from the saved run files (\texttt{stats.json}, \texttt{progress.log}, \texttt{fails.jsonl}) and the local certificate: the run finished and its last log line agrees with the statistics; $\mathrm{fail} = 0$ and the failure log is empty; leaves minus internal nodes equals $36$, the number of root boxes; $0 < T \le t_{\rm lower}(1-10^{-9}) < t_{\rm lower}$ in exact rationals, with the $t_{\rm lower}$ of the certificate; and the float weight tally is within $10^{-9}$ of $1$. It passes. This is a consistency check of the counters the run recorded; it does not re-derive any box or re-run any test, so it is not an independent check of the computation. The $32$ covered boxes lie around the symmetry images of the optimum that are compatible with the diagonal reduction. % SRC: arithmetic on stats.json; 110/NOTES.md Result (fails.jsonl empty, stderr empty, covers)

% =====================================================================
\section{Limitations and what an independent check needs}\label{sec:limits}

\begin{itemize}
\item \emph{Not independently checked.} As stated in Remark~\ref{rem:status}, this is an AI-assisted computer proof that only its author has examined. % SRC: 110/NOTES.md "Not yet independently reviewed"
\item \emph{No stored certificates.} The leaves of the tree are not stored (there are about $1.8\cdot10^8$); the proof can be re-checked only by re-running the search, about $16$ hours on $14$ workers. Every P2 discard was decided by the exact integer certificate; the independent rational re-derivation (audit) was additionally run on a random $0.02\%$ of them (about $1$ in $5\,000$). % SRC: 103/NOTES.md Limitations; stats.json audited / p2
\item \emph{Trust base} as in Section~\ref{sec:sound}; none of it formally verified. P1 at $n = 9$ uses the 64-bit integer kernel of experiment~110, whose agreement with a rational computation was tested only on samples.
\item \emph{Identification.} The equality $x^* = C_9$ is checked exactly in $\Q(\sqrt{65})$, but uses two values from a re-run of the certifier that are not stored (Remark~\ref{rem:ident}); the optimality and uniqueness statements do not depend on it.
\end{itemize}

An independent check would need at least the following.
\begin{enumerate}[label=(\roman*)]
\item \emph{Re-running the verifier}: running the published code from the repository on other hardware and obtaining zero failures. Node counts need not match exactly: the hot-started LP depends on the order in which a worker processes boxes, so the duals, and hence which boxes P2 discards, can vary between runs.
\item \emph{An independent implementation} of the box tests P0--P2 and of the cover test, ideally emitting per-leaf certificates that a small, separately written checker (or a proof assistant) verifies; and an independent check of the local certificate of Theorem~\ref{thm:local} at $n = 9$.
\item \emph{Review of the reductions}: Lemma~\ref{lem:norm} (stretching, the diagonal, $x$- and $y$-mirrors, labelling), the correctness of the outward rounding in P0, the corner argument in P1, the validity of the McCormick and area rows on a box in P2, and the deduction of Theorem~\ref{thm:main} in Section~\ref{sec:outline}.
\end{enumerate}

% =====================================================================
\section*{Data and code availability}
The code (branch and bound, drivers, tests, local certifier, the exact identification script \texttt{check\_closed\_form\_exact.py}, the run-accounting check \texttt{verify\_run.py}, the certificate dump \texttt{dump\_p2.py} and the clean-room checker \texttt{check\_p2.py}), the run logs and statistics, and the sampled $n = 9$ certificates are available in the public repository accompanying this paper, \repourl, with SHA-256 checksums of the proof inputs. % SRC: task instructions (repository placeholder)

\section*{Use of AI tools}
The author used Claude, an AI system developed by Anthropic, for computation (writing and running search, certification and analysis code) and for drafting text. All claims were checked against the computational records; the author takes responsibility for the content. AI systems are not authors of this work.

% =====================================================================
\appendix

\section{The P2 certificate}\label{app:p2}
% SRC: 110-n9-proof/bb2.py Problem2._build_lp, k_fill, k_matvec, lp_bound2, audit2, process; 103-global-proof/bb.py Problem.__init__ (tri, pairs, tw, tws, colscale)

This appendix is derived from \texttt{bb2.py} (experiment~110) and the parts of \texttt{bb.py} (experiment~103) that it inherits. Throughout, $n = 9$, $D = 2^{24}$, the $m = 84$ triples are $t = (i,j,k)$ with $i<j<k$, and the $72$ ordered pairs are $p = (a,b)$, $a \ne b$.

\subsection*{A.1 The box and the orientation pattern}
P2 is applied to a box $B = \prod_k [\ell_k/D, h_k/D]$ \emph{after} P0 and P1 have run on it. P0 changes the integer endpoints in place, only raising $\ell_k$ and lowering $h_k$, and returns ``empty'' as soon as some $\ell_k > h_k$; so on every box that reaches P2, $0 \le \ell_k \le h_k \le D$. Coordinates $k < 9$ are $x_k$, coordinates $9 + i$ are $y_i$.

P1 computes, for each triple $t$, the minimum $m_t$ and maximum $M_t$ of the integer $D^2\cdot 2\Delta_t$ over the $64$ corner tuples, and sets
\[
 \sigma_t = +1 \ \text{if } m_t \ge -T_2,\qquad \sigma_t = -1 \ \text{if not, and } M_t \le T_2,\qquad \sigma_t = 0 \ \text{otherwise},
\]
with $T_2 = \lfloor 2TD^2\rfloor$ (and discards the box if $\max(M_t,-m_t) \le T_2$ for some $t$). Since $2\Delta_t = x_jy_k - x_jy_i - x_iy_k - x_ky_j + x_ky_i + x_iy_j$ is multilinear in the six coordinates, $m_t \le D^2\,2\Delta_t(P) \le M_t$ for every $P \in B$.

\begin{lemma}\label{lem:orient}
If $\sigma_t \ne 0$ and $P \in B$ has $\lvert\Delta_t(P)\rvert > T$, then $\sigma_t\,2\Delta_t(P) > 2T$.
\end{lemma}
\begin{proof}
If $\sigma_t = +1$, then $D^2\,2\Delta_t(P) \ge m_t \ge -T_2 \ge -2TD^2$, so $2\Delta_t(P) < -2T$ is impossible, and $\lvert 2\Delta_t(P)\rvert > 2T$ forces $2\Delta_t(P) > 2T$. The case $\sigma_t = -1$ is symmetric.
\end{proof}

There is no case split on orientations. The pattern $\sigma$ is a deduction valid for every configuration in $B$ that matters (those with $f > T$); a triple with $\sigma_t = 0$ simply contributes no row. So the relaxation below contains every $P \in B \cap \Omega$ with $f(P) > T$, and nothing has to be covered by other cases.

\subsection*{A.2 Variables, rows and bounds}
The integer model has $91$ columns, in this order: $u_0 = \hat\tau$; $u_{1+i} = X_i$ and $u_{10+i} = Y_i$ ($i = 0,\dots,8$); $u_{19+p} = W_p$ for the $72$ pairs. For a configuration $P$ the \emph{intended values} are
\[
 X_i = Dx_i,\quad Y_i = Dy_i,\quad W_{ab} = D^2 x_a y_b,\quad \hat\tau = D^2 \min_{t:\ \sigma_t \ne 0} \sigma_t\, 2\Delta_t(P),
\]
so that $D^2\,2\Delta_t = L_t(W) := W_{ij} - W_{ik} + W_{jk} - W_{ji} + W_{ki} - W_{kj}$. The $465$ rows $G_\iota u \le H_\iota$ have integer coefficients; each is $D^2$ times a row in natural units. For a pair $p = (a,b)$ write $X^-_p = \ell_a$, $X^+_p = h_a$, $Y^-_p = \ell_{9+b}$, $Y^+_p = h_{9+b}$.
\begin{enumerate}[label=(R\arabic*)]
\item \emph{Area rows} ($168$): for each triple $t$, row $2t$ is $\hat\tau - L_t(W) \le 0$ and row $2t+1$ is $\hat\tau + L_t(W) \le 0$. Row $2t$ is \emph{switched on} iff $\sigma_t = +1$, row $2t+1$ iff $\sigma_t = -1$. The rows are present in the LP for all triples (fixed structure); a switched-off row has upper bound $+\infty$ in the float LP.
\item \emph{McCormick rows} ($4\cdot72 = 288$), for each pair $p = (a,b)$:
\[
\begin{aligned}
 -W_p + X^-_p Y_b + Y^-_p X_a &\le X^-_p Y^-_p, &\qquad -W_p + X^+_p Y_b + Y^+_p X_a &\le X^+_p Y^+_p,\\
 W_p - X^+_p Y_b - Y^-_p X_a &\le -X^+_p Y^-_p, & W_p - X^-_p Y_b - Y^+_p X_a &\le -X^-_p Y^+_p .
\end{aligned}
\]
They are $D^2$ times $(x_a - x^-)(y_b - y^-) \ge 0$, $(x^+ - x_a)(y^+ - y_b) \ge 0$, $(x^+ - x_a)(y_b - y^-) \ge 0$, $(x_a - x^-)(y^+ - y_b) \ge 0$ with $w = x_ay_b$, and hold at the intended values for \emph{every} $P \in B$. Their coefficients are rewritten for each box.
\item \emph{Ordering rows} ($8$): $D X_i - D X_{i+1} \le 0$, valid on $\Omega$.
\item \emph{Mean row} ($1$): $D\sum_i X_i \le D^2 n/2 = 9\cdot 2^{47}$, valid on $\Omega$.
\end{enumerate}
Column bounds used by the certificate (the $\hat\tau$ bound is not used): $X_i \in [\ell_i, h_i]$, $Y_i \in [\ell_{9+i}, h_{9+i}]$, and $W_{ab} \in [\underline W_{ab}, \overline W_{ab}]$, the minimum and maximum of the four integer products $\{\ell_a,h_a\}\cdot\{\ell_{9+b},h_{9+b}\}$; since $x_a y_b$ is bilinear on a rectangle, $D^2x_ay_b$ lies in this range for every $P \in B$. We write $[\underline u_c, \overline u_c]$ for the bound of column $c \ge 1$.

\subsection*{A.3 The multipliers}
HiGHS solves $\max \hat\tau$ over these rows in natural units (floating point). Let $\lambda = \max(-\text{row dual}, 0)$ componentwise. Then the code sets $\lambda_\iota = 0$ for every area row that is switched off (including both rows of every triple with $\sigma_t = 0$, and the wrong-orientation row of every determined triple). If $\lambda_{\max} = \max_\iota\lambda_\iota$ is not positive, P2 fails. Otherwise $S = 49 - \max(0, \lceil\log_2\lambda_{\max}\rceil)$ and $\mu_\iota = \lfloor \lambda_\iota 2^S\rfloor$; the code asserts $0 \le \mu_\iota < 2^{50}$. Nothing below depends on how $\mu$ was obtained: the solver, its tolerances, the scaling and the rounding only influence whether the test succeeds.

\subsection*{A.4 The exact test}
Put $r = G^{\mathsf T}\mu \in \mathbb{Z}^{91}$, $s = r_0$ (the $\hat\tau$ coefficient; only area rows have one, equal to $+1$), and
\[
 N = \mu^{\mathsf T} H + \sum_{c=1}^{90} \max\bigl(-r_c\,\underline u_c,\ -r_c\,\overline u_c\bigr) \in \mathbb{Z}.
\]
Let $2TD^2 = \mathfrak p/\mathfrak q$ in lowest terms. The box is discarded by P2 if and only if $s > 0$ and $N\mathfrak q \le s\mathfrak p$, that is, $N \le 2TD^2 s$; the comparison is in Python integers.

\begin{lemma}[P2 certificate]\label{lem:farkas}
Let $B$ be a box with integer endpoints in $[0,D]$ and orientation pattern $\sigma$ as in A.1, and let $\mu \in \mathbb{Z}_{\ge0}^{465}$ with $\mu_\iota = 0$ for every switched-off area row. If $s > 0$ and $N \le 2TD^2 s$, then no $P \in B \cap \Omega$ has $f(P) > T$.
\end{lemma}
\begin{proof}
Suppose $P \in B\cap\Omega$ and $f(P) > T$, and let $u$ be the intended values. Since $s = \sum_{\iota\ \text{area}} \mu_\iota > 0$ and $\mu$ vanishes on switched-off rows, some triple has $\sigma_t \ne 0$; so $\hat\tau$ is defined, and by Lemma~\ref{lem:orient} (every $\lvert\Delta_t(P)\rvert > T$) $\hat\tau > 2TD^2$. Every row with $\mu_\iota > 0$ holds at $u$: a switched-on area row of triple $t$ reads $\hat\tau \le \sigma_t D^2 2\Delta_t(P)$, true by definition of $\hat\tau$; McCormick rows hold since $P \in B$; ordering and mean rows hold since $P \in \Omega$. Multiplying by $\mu_\iota \ge 0$ and adding,
\[
 s\hat\tau + \sum_{c \ge 1} r_c u_c = \mu^{\mathsf T} G u \le \mu^{\mathsf T} H .
\]
Each $u_c$ ($c \ge 1$) lies in $[\underline u_c, \overline u_c]$, so $-r_c u_c \le \max(-r_c\underline u_c, -r_c\overline u_c)$. Hence $s\hat\tau \le N \le 2TD^2 s$, and since $s > 0$, $\hat\tau \le 2TD^2$, a contradiction.
\end{proof}

The lemma gives exactly what P2 is used for in (S2): $\min_{\sigma_t\ne0}\sigma_t 2\Delta_t \le 2T$ on $B\cap\Omega$, so no configuration with $f > T$ (configurations with $f = T$ exactly are not excluded, and need not be, since $T < \tstar$).

\subsection*{\texorpdfstring{A.5 Exactness of $r = G^{\mathsf T}\mu$}{A.5 Exactness of r}}
The products are computed in numba's 64-bit integers, which wrap silently on overflow; overflow is excluded as follows. Write $\mu = 2^{25}\mu^{\rm h} + \mu^{\rm l}$ with $0 \le \mu^{\rm h}, \mu^{\rm l} < 2^{25}$ (possible since $\mu < 2^{50}$), compute $r^{\rm h} = G^{\mathsf T}\mu^{\rm h}$ and $r^{\rm l} = G^{\mathsf T}\mu^{\rm l}$ in int64, and recombine $r = 2^{25}r^{\rm h} + r^{\rm l}$ in Python integers; $N$ is formed in Python integers. Every entry of $G$ has absolute value at most $D = 2^{24}$: area and $W$-coefficients are $\pm1$, ordering and mean coefficients $\pm D$, and McCormick coefficients are endpoints in $[0,D]$ (A.1; \texttt{bb2.py} does not re-assert this at run time, \texttt{bb.py} did). So each product is below $2^{24}\cdot2^{25} = 2^{49}$ in absolute value. Column counts: $\hat\tau$ has $2m = 168$ nonzeros; $W_{ab}$ has $14 + 4 = 18$ (seven triples contain $\{a,b\}$, two rows each, plus four McCormick rows); $X_a$ at most $32 + 2 + 1 = 35$; $Y_b$ has $32$. Hence every partial sum is below $168\cdot2^{49} < 2^{57}$. The code computes the maximum column count when the model is built and asserts $\text{count}\cdot D\cdot2^{25} < 2^{62}$. The McCormick right-hand sides and $W$ bounds are products of two endpoints, at most $2^{48}$, and the mean right-hand side is $9\cdot2^{47}$.

\subsection*{A.6 What the lemma does not cover}
Everything that decides a P2 discard is covered: the code discards only through the test of A.4, applied to the same box (after P0) whose endpoints were written into $G$, $H$ and the bounds immediately before. The following parts of the code are outside the lemma and cannot discard a box: the early return when the float LP value exceeds $2T(1+10^{-7})+10^{-12}$; non-optimal solver status; $\lambda_{\max} \le 0$ or $s \le 0$; LP skips. The audit (a rational re-derivation of a random $0.02\%$ of certificates, which also rejects any nonzero multiplier on a switched-off row) is an additional check and not part of the lemma. The lemma does rely on two facts about the implementation that we checked by reading the code, not by proof: that the zeroing \texttt{lam[:2*m][\textasciitilde act] = 0} acts on a view of \texttt{lam} (it does: a basic slice of a NumPy array is a view), and that numba's int64 arithmetic is exact in the absence of overflow.

% =====================================================================
\section{The local certificate and its basin radius}\label{app:basin}
% SRC: 44-interval-certificate/certify_interval.py (build_param, contraction, hess_l1, precond_inverse, tie_row, certify l.433-547); 103-global-proof/results/local_certs.json key "9"; re-run of certify(9) on 102-laws-small-n optima_small.json (values Y, Z0, Z2, L, ||C||, beta, y bounds, parametrisation)

This appendix is derived from \texttt{certify\_interval.py} (experiment~44), run unchanged on the $n = 9$ optimum by \texttt{get\_certs.py} (experiment~103); the numbers are those of the stored record (\texttt{local\_certs.json}, key \texttt{"9"}) and, for intermediate quantities not stored there, of a re-run of the same deterministic code, which reproduces the stored record exactly. Norms: $\|v\|_\infty = \max_i\lvert v_i\rvert$ for vectors and $\|A\|_\infty = \max_i\sum_k\lvert A_{ik}\rvert$ for matrices. $A_T$ denotes the \emph{signed} area of a triple $T = (i,j,k)$, $i<j<k$ (not doubled).

\subsection*{B.1 The reduced tie system}
The float stage finds $11$ ties (relative tolerance $10^{-6}$) and $8$ wall contacts (tolerance $10^{-8}$), all positively loaded in the KKT stress LP, no rattlers, and the symmetry group $\{\mathrm{id}, (x,y)\mapsto(y,x)\}$. With the labels of $\tilde x$ (the order of $C_9$ in Section~\ref{sec:intro}), the exact equivariant, wall-respecting affine parametrisation built by the certifier is, with $u \in \R^5$,
\begin{equation}\label{eq:param}
\begin{gathered}
 p_0 = (-u_0, 1),\ p_1 = (1, -u_0),\ p_2 = (1, u_1),\ p_3 = (u_1, 1),\ p_4 = (0, u_2),\\
 p_5 = (u_2, 0),\ p_6 = (-u_3, 0),\ p_7 = (0, -u_3),\ p_8 = (-u_4, -u_4).
\end{gathered}
\end{equation}
Every coordinate is a constant or $\pm$ one parameter, so the maximal row sum of the coefficient matrix is $b_{\max} = 1$. The $11$ ties form $6$ orbits under the reflection ($5$ of size $2$, one of size $1$). With one representative $T_k$ and its sign $\sigma_k$ per orbit, the unknowns $w = (u, t) \in \R^6$ and
\[
 F_k(w) = \sigma_k A_{T_k}(p(u)) - t, \qquad k = 1,\dots,6 .
\]
The Jacobian has numerical rank $6$ (no orbit is dropped): \emph{first-order mode}. Newton's method, with exact residuals and iterates rounded to the grid $2^{-110}$, gives $\tilde w = (\tilde u, \tilde t)$ with $\|F(\tilde w)\|_\infty \approx 8.5\cdot10^{-33}$; $\tilde x = p(\tilde u)$ is stored as \texttt{points\_exact} (dyadic rationals with denominators up to $2^{108}$).

\subsection*{B.2 Existence and uniqueness}
Since $F$ is quadratic, $DF_k(w) - DF_k(\tilde w) = (\sigma_k H_k(u - \tilde u), 0)$ with $H_k$ the constant Hessian of $u \mapsto A_{T_k}(p(u))$, so $\|DF_k(w) - DF_k(\tilde w)\|_1 \le L_k\|w - \tilde w\|_\infty$ with $L_k = \sum_{l,l'}\lvert(H_k)_{ll'}\rvert$; for $n = 9$, $L = (2,2,2,2,2,3)$. Let $R$ be the float inverse of $DF(\tilde w)$ rounded to dyadic rationals ($50$-bit integer matrix over a power of two). The code computes, in exact rational arithmetic,
\[
 Y = \|RF(\tilde w)\|_\infty,\qquad Z_0 = \|I - R\,DF(\tilde w)\|_\infty,\qquad Z_2 = \max_i \sum_k \lvert R_{ik}\rvert L_k .
\]
\begin{lemma}\label{lem:krawczyk}
If $r > 0$ satisfies $Y + Z_0 r + Z_2r^2/2 \le r$ and $Z_0 + Z_2 r < 1$, then $F$ has exactly one zero $w^*$ in the closed ball $\bar B = \{\|w - \tilde w\|_\infty \le r\}$.
\end{lemma}
\begin{proof}
Let $\Phi(w) = w - RF(w)$. For $w \in \bar B$, $D\Phi(w) = I - R\,DF(\tilde w) - R(DF(w) - DF(\tilde w))$ and $\|R(DF(w) - DF(\tilde w))\|_\infty \le Z_2\|w-\tilde w\|_\infty$ (row $i$: $\sum_k\lvert R_{ik}\rvert L_k\|w - \tilde w\|_\infty$), so $\|D\Phi(w)\|_\infty \le Z_0 + Z_2\|w - \tilde w\|_\infty$. Then $\|\Phi(w) - \tilde w\|_\infty \le \|\Phi(\tilde w) - \tilde w\|_\infty + \int_0^1 \|D\Phi(\tilde w + \theta(w-\tilde w))\|_\infty\, r\,d\theta \le Y + Z_0r + Z_2r^2/2 \le r$, so $\Phi(\bar B) \subset \bar B$, and $\Phi$ is a contraction on $\bar B$ with constant $Z_0 + Z_2r < 1$. It has exactly one fixed point. Since $Z_0 < 1$, $R\,DF(\tilde w)$ is invertible, hence so is $R$, and the fixed points of $\Phi$ are exactly the zeros of $F$.
\end{proof}
The code tries $r = \max(2Y/(1-Z_0), 10^{-40})$ and multiplies by $4$ until both inequalities hold. For $n = 9$: $Y = 2.3247\cdot10^{-32}$, $Z_0 = 6.78\cdot10^{-15}$, $Z_2 = 15.31$, and the first trial $r = 2Y/(1-Z_0) = 4.6494\cdot10^{-32}$ succeeds (the first inequality reduces to $Z_2 r \le 1 - Z_0$). Write $w^* = (u^*,\tstar)$ and $x^* = p(u^*)$. Every coordinate of $x^*$ is within $\rho = b_{\max} r = r$ of $\tilde x$, and $\tstar \in [\tilde t - r, \tilde t + r] \subset [t_{\rm lower}, t_{\rm upper}]$ (the endpoints are rounded outward to $32$ decimals).

\subsection*{\texorpdfstring{B.3 The exact active set at $x^*$}{B.3 The exact active set at x*}}
At $x^*$ the six solved equations and their images under the reflection give $11$ triangles with $\sigma_TA_T(x^*) = \tstar$ exactly ($\sigma_T$ the sign at $\tilde x$; it is the sign at $x^*$ because $\lvert A_T\rvert \approx \tstar$ and $\rho$ is tiny), and $x^*$ has the $8$ wall contacts of \eqref{eq:param} exactly. (The check that all other triangles exceed $\tstar$ at $x^*$ is needed for $f(x^*) = \tstar$, not for the basin inequality; for $n = 9$ the certifier does it in double precision with margin $\ge 10^{-6}$, and it is superseded by the exact check of Lemma~\ref{lem:ident}.)

\subsection*{B.4 The stress matrix and the multipliers}
Let $N = 2\cdot9 + 1 = 19$. Columns $2j, 2j+1$ of the $19\times19$ matrix $M(x)$ belong to $x_j, y_j$, column $18$ to an auxiliary variable $\tau$. Rows: for each of the $11$ exact ties $T$, the row $[\sigma_T\nabla A_T(x),\, -1]$, where $\partial A_T/\partial x_i = (y_j - y_k)/2$, $\partial A_T/\partial y_i = (x_k - x_j)/2$ and cyclically; for each of the $8$ wall contacts $(j,c)$, the row $[n_w e_{2j+c},\, 0]$ with $n_w = +1$ on a side $\{\cdot = 0\}$ and $n_w = -1$ on a side $\{\cdot = 1\}$. Since there are exactly $19$ exact constraints, all are used (no vertex selection). A tie row has $6$ coordinate entries, each half a difference of two coordinates, so each moves by at most $\rho$ between $\tilde x$ and $x^*$:
\[
 \|M(x^*) - M(\tilde x)\|_\infty \le 6\rho .
\]
Let $C$ be the float inverse of $M(\tilde x)$, rounded to a dyadic matrix, and
\[
 \beta = \|I - C M(\tilde x)\|_\infty + 6\rho\,\|C\|_\infty \qquad(\text{exact}).
\]
\begin{lemma}\label{lem:stress}
If $\beta < 1$, then $M^* = M(x^*)$ is invertible, $\|(M^*)^{-1}\|_\infty \le \bar m := \|C\|_\infty/(1-\beta)$, and every entry satisfies $\lvert ((M^*)^{-1} - C)_{ik}\rvert \le \varepsilon := \|C\|_\infty\beta/(1-\beta)$. Consequently $y := -(\text{row } 18 \text{ of } (M^*)^{-1})$ satisfies $y^{\mathsf T}M^* = -e_{18}^{\mathsf T}$ and $y_k \ge y_k^- := -C_{18,k} - \varepsilon$.
\end{lemma}
\begin{proof}
$E = I - CM^*$ has $\|E\|_\infty \le \|I - CM(\tilde x)\|_\infty + \|C\|_\infty\|M(\tilde x) - M^*\|_\infty \le \beta < 1$, so $CM^* = I - E$ is invertible with $\|(CM^*)^{-1}\|_\infty \le 1/(1-\beta)$; hence $C$ and $M^*$ are invertible and $(M^*)^{-1} = (CM^*)^{-1}C$. Further $(M^*)^{-1} - C = ((CM^*)^{-1} - I)C = (CM^*)^{-1}EC$, of norm at most $\beta\|C\|_\infty/(1-\beta)$, which bounds each entry.
\end{proof}
For $n = 9$: $\|C\|_\infty = 46.0492$, $\|I - CM(\tilde x)\|_\infty = 3.70\cdot10^{-14}$, $6\rho\|C\|_\infty = 1.3\cdot10^{-29}$, $\beta = 3.70\cdot10^{-14}$, $\varepsilon = 1.70\cdot10^{-12}$, $\bar m = 46.04921456503519$. The lower bounds $y_k^-$ are $0.0464$--$0.1128$ on the tie rows and $0.0217$--$0.0332$ on the wall rows; $y_{\min} := \min_k y_k^- = 0.0217100466552658$.

\subsection*{B.5 Uniform decrease and the basin radius}
Put $c = y_{\min}/\bar m$ ($= 4.71453\cdot10^{-4}$ for $n = 9$).
\begin{lemma}[uniform decrease]\label{lem:decrease}
Let $d \in \R^{18}$ with $x^* + d \in Q^9$ and $\delta = \|d\|_\infty > 0$. Then some exact tie $T$ has $\sigma_T\nabla A_T(x^*)\cdot d \le -c\,\delta$.
\end{lemma}
\begin{proof}
Let $\tau = \min_T \sigma_T\nabla A_T(x^*)\cdot d$ over the $11$ ties and $s = M^*(d,\tau)$. A tie row gives $s_T = \sigma_T\nabla A_T(x^*)\cdot d - \tau \ge 0$. A wall row gives $s_w = n_w d_{2j+c} \ge 0$, because the coordinate is $0$ (resp.\ $1$) at $x^*$ and stays in $[0,1]$. So $s \ge 0$, and by Lemma~\ref{lem:stress}
\[
 -\tau = y^{\mathsf T}M^*(d,\tau) = y\cdot s \ge y_{\min}\sum_k s_k \ge y_{\min}\|s\|_\infty .
\]
Also $\delta \le \|(d,\tau)\|_\infty = \|(M^*)^{-1}s\|_\infty \le \bar m\|s\|_\infty$. Hence $-\tau \ge (y_{\min}/\bar m)\,\delta = c\delta$, and the minimising tie $T$ has $\sigma_T\nabla A_T(x^*)\cdot d = \tau \le -c\delta$.
\end{proof}
This is the inequality that turns strictly positive multipliers into a decrease: positivity of every $y_k$ makes $y\cdot s$ control $\|s\|_\infty$, and invertibility of $M^*$ makes $\|s\|_\infty$ control $\delta$.

\begin{theorem}\label{thm:basin}
Let $\rho_b = \min(c/4, \tilde t/10)$. Every $x \in Q^9$ with $0 < \|x - x^*\|_\infty < \rho_b$ has $f(x) < \tstar$. For $n = 9$, $\rho_b = c/4$, stored as the double $\rho_b' = 1.1786328420762071\cdot10^{-4}$.
\end{theorem}
\begin{proof}
Let $d = x - x^*$, $0 < \delta < c/4$, and take $T = (i,j,k)$ from Lemma~\ref{lem:decrease}. Since $A_T$ is quadratic,
\[
 A_T(x) = A_T(x^*) + \nabla A_T(x^*)\cdot d + \tfrac12\bigl[(d_j - d_i)\times(d_k - d_i)\bigr],
\]
with $\times$ the planar cross product; each component of $d_j - d_i$, $d_k - d_i$ has absolute value $\le 2\delta$, so the last term is at most $\tfrac12(4\delta^2 + 4\delta^2) = 4\delta^2$ in absolute value. Hence
\[
 \sigma_TA_T(x) \le \tstar - c\delta + 4\delta^2 = \tstar - \delta(c - 4\delta) < \tstar .
\]
On the other side, $\|\nabla A_T\|_1 \le 6\cdot\frac12 = 3$ (coordinates in $[0,1]$), so $\sigma_TA_T(x) \ge \tstar - 3\delta - 4\delta^2 > 0$ because $\delta \le \tilde t/10$ and $\tilde t$ is within $10^{-31}$ of $\tstar$. Thus $\lvert A_T(x)\rvert < \tstar$ and $f(x) < \tstar$.
\end{proof}
Since $y_{\min}$ and $\bar m$ are a lower and an upper bound, $c$ is a lower bound for the true $y$-ratio, which is all the proof uses. For $n = 9$ the cap $\tilde t/10 \approx 5.5\cdot10^{-3}$ is not binding.

\subsection*{B.6 How the branch and bound uses the radius}
\texttt{bb.py} (\texttt{load\_covers}) reads the stored doubles $\rho_b'$ and $\rho' = $ \texttt{box\_rho} $= 4.649442470662108\cdot10^{-32}$ as exact rationals and sets
\[
 R = \rho_b'(1 - 10^{-6}) - \rho' - 2^{-60} \approx 1.1786316634\cdot10^{-4}.
\]
\emph{Claim:} $R + \rho < \rho_b$, where $\rho$ is the exact point-displacement bound. Indeed $\rho - \rho' \le 2^{-53}\rho < 2^{-60}$ (round to nearest), so $R + \rho < \rho_b'(1-10^{-6})$; and $\rho_b' \le \rho_b(1 + 2^{-53})$, so $\rho_b'(1 - 10^{-6}) < \rho_b$.
The cover test (\texttt{covered}) is applied to a box $B$ (after P0) of width at most $D/1000$. Floating point proposes a symmetry $g$ (one of the eight maps $c + M(p - c)$, $c = (\frac12,\frac12)$, applied to $\tilde x$ in exact rationals) and an assignment $i \mapsto j(i)$ of box points to points of $g(\tilde x)$, which must be injective (checked on integers). The box is covered iff, in exact rationals, for every $i$ and both coordinates, both endpoints $\ell/D$ and $h/D$ are within $R$ of the coordinate of $g(\tilde x)_{j(i)}$. As the sup-ball is a product of intervals, every $P \in B$ then has $\|P - \pi(g(\tilde x))\|_\infty \le R$ with $\pi$ the relabelling, so $\|g^{-1}(\pi^{-1}P) - x^*\|_\infty \le R + \rho < \rho_b$ (symmetries of the square are sup-norm isometries), and Theorem~\ref{thm:basin} applies to $g^{-1}(\pi^{-1}P)$.

\subsection*{B.7 Points not justified by the code alone}
(i) Two values used in Lemma~\ref{lem:ident}, the 32-digit display of $\tilde t$ and the parametrisation \eqref{eq:param}, come from a re-run of the certifier and are not stored (Remark~\ref{rem:ident}). (ii) The certifier's check that non-tied triangles exceed $\tstar$ uses double precision for margins above $10^{-6}$; for $n = 9$ this is replaced by the exact check of Lemma~\ref{lem:ident}. (iii) The cap uses $\tilde t$ rather than $\tstar$; the proof of Theorem~\ref{thm:basin} shows this is harmless. We found no step of B.2--B.6 that the code does not justify.

% =====================================================================
\begin{thebibliography}{99}

\bibitem{Chen17} L.~Chen, Y.~Xu and Z.~Zeng, Searching approximate global optimal Heilbronn configurations of nine points in the unit square via GPGPU computing, \emph{J. Global Optim.} \textbf{68} (2017), 147--167, doi:10.1007/s10898-016-0453-1.
\bibitem{CY} F.~Comellas and J.~L.~A.~Yebra, New lower bounds for Heilbronn numbers, \emph{Electron. J. Combin.} \textbf{9} (2002), \#R6, doi:10.37236/1623.
\bibitem{DZ} L.~Dehbi and Z.~Zeng, Heilbronn's problem of eight points in the square, \emph{J. Syst. Sci. Complex.} \textbf{35} (2022), 2452--2480, doi:10.1007/s11424-022-1220-7.
\bibitem{DYZ} A.~W.~M.~Dress, L.~Yang and Z.~Zeng, Heilbronn problem for six points in a planar convex body, in: D.-Z.~Du and P.~M.~Pardalos (eds.), \emph{Minimax and Applications}, Kluwer, 1995, 173--190, doi:10.1007/978-1-4613-3557-3\_13.
\bibitem{MMK} A.~Monji, A.~Modir and B.~Kocuk, Solving the Heilbronn triangle problem using global optimization methods, arXiv:2512.14505v1, 2025.
\bibitem{Stead} T.~Stead (maintainer), The Heilbronn Problem (record tables, exact verification and attributions), \url{https://math.tejstead.com/heilbronn}, accessed 12 November 2026.
\bibitem{SM} N.~Sudermann-Merx, From global optimization to exact coordinates: Heilbronn's triangle problem on the unit square using mixed-integer nonlinear optimization, arXiv:2603.11107v3, 27 October 2026 (v1: 11 March 2026).
\bibitem{YZZ92} L.~Yang, J.~Z.~Zhang and Z.~B.~Zeng, A conjecture on the first several Heilbronn numbers and a computation, \emph{Chinese Ann. Math. Ser. A} \textbf{13} (1992), 503--515.
\bibitem{YZ80} L.~Yang and J.~Zhang, The six-point problem in the square, in: \emph{Mathematics Lectures (Part II)}, Sichuan People's Publishing House, 1980, 159--175 (in Chinese).
\bibitem{YZZ} L.~Yang, J.~Zhang and Z.~Zeng, Heilbronn problem for five points, Technical Report IC/91/252, International Centre for Theoretical Physics, Trieste, 1991.
\bibitem{ZC} Z.~Zeng and L.~Chen, On the Heilbronn optimal configuration of seven points in the square, in: \emph{Automated Deduction in Geometry} (ADG 2008), Lecture Notes in Comput. Sci. \textbf{6301}, Springer, 2011, 196--224, doi:10.1007/978-3-642-21046-4\_11.

\bibitem{Pan10} C.~Panorias, An exact computer-assisted proof for ten points in Heilbronn's triangle problem, preprint, version~1.1, Zenodo, 13 November 2026, \url{https://zenodo.org/records/24186293}.
\end{thebibliography}

\end{document}
